\documentclass[11pt]{article}
\usepackage[a4paper,margin=27mm]{geometry}
\usepackage{amsmath,amssymb,amsthm,booktabs,hyperref}
\newtheorem{theorem}{Theorem}
\newtheorem{lemma}[theorem]{Lemma}
\newtheorem{corollary}[theorem]{Corollary}
\newcommand{\F}{\mathbb F}
\newcommand{\PG}{\operatorname{PGL}}
\newcommand{\PS}{\operatorname{PSL}}
\newcommand{\AG}{\operatorname{AGL}}
\title{Subfield planes in projective rank two geometries}
\author{Human Compression project}
\date{}
\begin{document}
\maketitle
\begin{abstract}
A plane over a subfield supplies a locally two-transitive coset geometry
without requiring the ambient field to have degree two over that subfield.
We give the construction and check its intersection, connectivity and
primitivity directly. For $\PS(2,512)$ the two parabolics have orders $504$
and $448$, with intersection of order $56$. This contradicts the exhaustion
in the fixed Part II classification and the corresponding restriction in
its separately catalogued appendix. It is one defect appearing in two
documents, not two independent discoveries. No complete corrected
classification is asserted.
\end{abstract}

\section{Scalar subspaces and local actions}
We write $V\rtimes K^*$ for the affine transformations $x\mapsto ax+v$,
where $K$ is a finite field, $V$ is a vector space over $K$, $a\in K^*$
and $v\in V$. Its multiplication is
$(v,a)(w,b)=(v+aw,ab)$. All actions on cosets below are allowed to have
a kernel; this is also the convention in the source classification.

\begin{lemma}\label{scalar}
Let $W\leq V$ be $K$-subspaces, and set $B=W\rtimes K^*$.
The subgroups between $B$ and $V\rtimes K^*$ are precisely
$T\rtimes K^*$ with $W\leq T\leq V$. Such a subgroup covers $B$ if
and only if $\dim_K(T/W)=1$. In this case its action on the cosets of
$B$ induces $\AG(1,K)$ and is two-transitive, with kernel $W$.
\end{lemma}
\begin{proof}
If $D$ contains $B$ and $(v,a)\in D$, multiplying by $(0,a^{-1})\in B$
shows that $(v,1)\in D$. Thus $D=(D\cap V)\rtimes K^*$.
The additive group $D\cap V$ is invariant under all nonzero scalars,
so it is a $K$-subspace containing $W$. This proves the correspondence
and the assertion about covers. Since $W$ is normal in $T\rtimes K^*$,
the induced coset action is the affine action of
$(T/W)\rtimes K^*$ on $T/W$. If this quotient has dimension one,
translations are regular and $K^*$ is regular on its nonzero vectors.
The action is therefore two-transitive and faithful on the quotient.
Its kernel before passage to the quotient is exactly $W$.
\end{proof}

The distinction between an additive plane and a subfield is essential.
For $F/K$ of prime degree $m>2$, there is no intermediate subfield of
degree two. Nevertheless there are two-dimensional $K$-subspaces of $F$.
Lemma~\ref{scalar} uses only a subspace.

\section{The projective construction}
Let $s=2^n\geq4$, let $m$ be prime, and put $K=\F_s\subset F=\F_{s^m}$.
In characteristic two every nonzero determinant has a square root, so
$\PG(2,F)=\PS(2,F)$. We use fractional linear transformations as literal
elements of this group. In particular the transformation $x\mapsto ax+b$
is represented by $\left(\begin{smallmatrix}a&b\\0&1\end{smallmatrix}\right)$.
Write $U_T=\{x\mapsto x+t:t\in T\}$ and
$C=\{x\mapsto ax:a\in K^*\}$. The retained classical subfield theorem
states that the natural $\PS(2,K)$ is maximal in $\PS(2,F)$ when $m$
is prime. This is a subgroup theorem used by the source, not its geometry
classification.

\begin{theorem}\label{main}
For every $K$-plane $T\subseteq F$ containing $K$, set
\[
 G=\PS(2,F),\qquad H=\PS(2,K),\qquad Q=U_T C,\qquad B=U_K C.
\]
Then $H\cap Q=B$, and the right-coset incidence system with parabolics
$H,Q$ is a thick, connected, flag-transitive rank two geometry. It is
residually weakly primitive and locally two-transitive. Its two local
degrees are $s+1,s$, and
\[
 |H|=s(s^2-1),\qquad |Q|=s^2(s-1),\qquad |B|=s(s-1).
\]
These assertions hold for every prime $m\geq2$.
\end{theorem}
\begin{proof}
The set $Q$ is a subgroup because $T$ is stable under multiplication by
$K$. Its elements have the unique affine expressions $x\mapsto ax+t$
with $a\in K^*$ and $t\in T$. Such an element belongs to $H$ precisely
when $t\in K$: a projective transformation over $K$ fixing infinity is
affine with both coefficients in $K$, and its normalized bottom right
entry is one. Thus the intersection is exactly $B$. The displayed
orders follow from these unique expressions and the classical order of
$\PS(2,K)$.

The action of $H$ on $H/B$ is its natural action on $\mathbb P^1(K)$.
It is two-transitive: the stabilizer of infinity contains all translations
of $K$, which are transitive on the other points. The action of $Q$ on
$Q/B$ is two-transitive by Lemma~\ref{scalar}, since $\dim_K(T/K)=1$.
The local degrees are consequently $s+1$ and $s$. In particular both
intersections are maximal in their parabolics. The second local action
has the nontrivial kernel $U_K$; faithfulness of local actions is not a
hypothesis of the source problem.

For completeness, take the elements of the incidence system to be the
right cosets of $H$ and $Q$, of distinct types, incident when they
intersect. Each incident pair contains an element $g$ and is therefore
the translate of $\{H,Q\}$ by $g$. This proves flag transitivity.
Its neighbours at the base vertices are the coset actions just computed,
so every vertex lies on an edge and both residues are thick.
Choose $t\in T\setminus K$. The translation by $t$ lies in $Q\setminus H$.
The retained maximality of $H$ therefore gives $\langle H,Q\rangle=G$.
Paths obtained by multiplying alternately by elements of $H$ and $Q$
show that the incidence graph is connected. In rank two this is the only
connectedness requirement for residues of rank at least two.

Finally $G$ acts primitively on $G/H$, since $H$ is maximal, and both
rank one residue actions are primitive, since they are two-transitive.
These are exactly the weak primitivity and residual primitivity
conditions in rank two. The whole action is faithful because its kernel
is a normal subgroup of the simple group $G$ contained in the proper
subgroup $H$.
\end{proof}

There are $(s^{m-1}-1)/(s-1)$ choices of $T$ for the fixed subfield and
fixed Borel $B$: they correspond to the lines of the $(m-1)$-dimensional
space $F/K$. This counts subspaces, not isomorphism classes of geometries.
For $m=2$ there is only $T=F$, recovering the known quadratic family.
For $m>2$ every such plane is a proper additive subspace of $F$.

\section{The two fixed sources}
The first source is Buekenhout, De Saedeleer and Leemans,
\emph{On the rank two geometries of the groups $\PS(2,q)$: part II},
Ars Mathematica Contemporanea 6 (2013), 365--388 \cite{partii}.
Its Theorem 1.1 asserts exhaustion by Tables 1--6. On printed page 382,
Lemma 5.11 restricts the affine overgroup $E_{2^{2n}}:(2^n-1)$ to $m=2$;
Proposition 5.12 and Table 4 on page 367 use that restriction in the
classification with subfield parabolic $\PS(2,2^n)$.

\begin{corollary}
The exhaustion asserted by Theorem 1.1, Proposition 5.12 and Table 4,
and the $m=2$ restriction in Lemma 5.11 of this fixed edition, are false.
\end{corollary}
\begin{proof}
Apply Theorem~\ref{main} with $s=8$, $m=3$. Then $q=512$, $n=3$,
and the parabolics have orders $504,448$ with intersection $56$.
The first row of Table 4 requires $m=2$, and every other row requires
$n=1$ or $n=2$. Tables 1--3 concern the maximal parabolics
$A_5,A_4,S_4$, whose orders are different from both $504$ and $448$.
Table 5 requires odd characteristic and Table 6 requires an odd square
field. Thus none of the tables contains this geometry, even after
interchanging its two types. Its local degrees $9,8$ and the verified
maximality and intersection satisfy the stated hypotheses. The same
group $Q$ is the forbidden affine overgroup in Lemma 5.11.
\end{proof}

The separately indexed appendix \cite{appendix} is a fixed 21-page
document with the placeholder heading \emph{Ars Mathematica Contemporanea
x (xxxx) 1--x}; its numbering differs from the published paper.
On page 12, the proof headed ``Proposition 22'', general case, subcase 2,
again admits this affine parabolic only for $m=2$. The construction above
contradicts that restriction and the resulting purported exhaustion.
The proof of Lemma 21 on page 11 actually lists the affine abstract
group without that restriction. We do not attribute a false restriction
to that passage, nor infer the full statement of its missing lemma from
the heading alone.

This is one omitted family shared by two catalogue entries. Other
classification cases and numerical counts in those sources have not
thereby been recovered or refuted. In particular no corrected global
class count or full automorphism group is claimed. Both complete source
texts were read to establish this boundary; only the indicated assertions
are contradicted here.

\section{Attribution and finite verification}
The quadratic example and the analogous $A_5$ families are already present
in the source. The general scalar-subspace correspondence is elementary;
the present application uses the same mechanism as the project's earlier
manuscript \emph{Scalar subspaces in rank-two Suzuki geometries}.
Here its role is to remove an unjustified extension-degree restriction.
The natural projective action, the order formula, simplicity and subfield
maximality are retained classical facts. This note neither claims
historical priority nor asserts that no later correction exists.

The accompanying program independently represents $\F_{512}$ as a binary
polynomial field, constructs the actual projective subgroups $H,Q,B$,
and checks their orders, intersections and the two local actions.
It also checks the quadratic boundary and the known $s=4$ specialization.
These finite controls supplement the proof; neither the uniform argument
nor the retained subfield maximality theorem is inferred from a finite
sample. No Lean verification or open-conjecture credit is asserted.

\begin{thebibliography}{9}
\bibitem{partii} F. Buekenhout, J. De Saedeleer and D. Leemans,
On the rank two geometries of the groups $\PS(2,q)$: part II,
\emph{Ars Math. Contemp.} \textbf{6} (2013), 365--388.
\href{https://amc-journal.eu/index.php/amc/article/download/181/241}{Fixed source}.
\bibitem{appendix} F. Buekenhout, J. De Saedeleer and D. Leemans,
On the rank two geometries of the groups $\PS(2,q)$: part II (APPENDIX),
the separately retained 21-page edition.
\href{https://amc-journal.eu/index.php/amc/article/download/181/1862}{Fixed appendix}.
\end{thebibliography}
\end{document}
